\section{Failure modes}\label{sec:failure-modes}

This section catalogues ten specific dead ends rejected during the
search programme. Each failure mode is paired with a regression fixture
in \texttt{tests/fixtures/} (or, when no falsifying graph is needed, a
combinatorial witness) so the rejection is reproducible. We include
this material in the main body (rather than an appendix) because the
\emph{retirement} of the universal ear lemma (\Cref{cor:universal-false})
and of single-scalar candidates is as informative for future work as
the theorems that survived.

\subsection{F1. The universal ear-deletion lemma}\label{sec:fm:universal}

The naive statement ``every simplicial degree-$2$ ear $v$ of every
$2$-tree on $n \ge 4$ vertices satisfies $\dminus(v) \ge 17/16$'' is
false. Two witnesses: the asymptotic $\dminus_{\infty}(\mathrm{BT}) \approx
1.0353 < 17/16$ on the $\BTkt{k}{2}$ tail family
(\Cref{thm:BT-asymptotic}; at $k = 50$ the value is already
$\dminus \approx 1.0575$), and a seed-$149$ random $2$-tree on
$n = 100$ vertices with $\dminus \approx 1.0610$. The latter graph is
preserved as the permanent regression fixture
\texttt{tests/fixtures/two\_tree\_universal\_counterexamples.json}; the
former is the structured proof of \Cref{cor:universal-false}. The
correct reformulation is the \emph{existential} ear-selection lemma
$\Lprime$ (\Cref{thm:Lprime}).

\subsection{F2. The $\|\earw\|^{2} = 4$ normalisation bug}\label{sec:fm:wnorm}

The supporting-edge vector $\earw = e_{a} + e_{b}$ has
$\|\earw\|^{2} = 2$, not $4$: the two standard basis vectors are
orthogonal, $e_{a}^{\top} e_{b} = 0$, regardless of whether
$\{a, b\}$ is an edge of $\Hgraph$. The mistaken value $4$ was
conflated with the trace identity $\dplus(v) + \dminus(v) = 2\deg_{\graph}(v) = 4$,
which is a different (and correct) derivation
($\tr(\Adj(\graph)^{2}) - \tr(\Adj(\Hgraph)^{2}) = 2 \deg_{\graph}(v)$).
Consequences of the corrected normalisation
$\sum_{i} \earc_{i}^{2} = \|\earw\|^{2} = 2$: the walk-moment identity
\Cref{lem:walk-moments} takes its stated form, the Cauchy--Schwarz pair
\Cref{lem:CS,lem:cs-two-sided} have the stated factors, and the
threshold $\threshold = 0.4122$ is no longer derivable from a
``natural'' rescaling of $17/16$. Regression fixture:
\texttt{tests/fixtures/w\_norm\_squared\_is\_2.json}.

\subsection{F3. Single-scalar selector thresholds}\label{sec:fm:single-scalar}

The natural single-scalar thresholds — $\Wm(\maxdeg) \ge 17/16$ (the
$\|\earw\|^{2} = 4$ era) and $\Wm(\maxdeg) \ge 17/32$ (after the v9
$\|\earw\|^{2} = 2$ renormalisation) — are both empirically false.
Witnesses: $L_{5}$ endpoint has $\Wm \approx 0.515$, $L_{6}$ has
$\Wm \approx 0.380$, and $L_{12}$ has $\Wm \approx 0.502$; all are
below $17/32 = 0.53125$. An empirically-fitted
threshold $T < 17/32$ does survive on a finite corpus, but the
worst-case bad ears from the $\BTkt{k}{2}$ tail family have
$\Wm \to 0$ as $k \to \infty$ while the max-degsum-ear infimum of
$\Wm$ also drops, narrowing the window asymptotically. The candidate
$\Iv = \Wm + (\Monem)^{2}/\Mtwom$ at $\threshold = 0.4122$
(\Cref{def:joint-invariant}) is a genuine joint invariant, not a
single-scalar rebrand. Regression test:
\texttt{tests/test\_lprime\_subfamilies.py}.

\subsection{F4. The naive sine-basis half-line boundary density}\label{sec:fm:sine-basis}

The naive guess for the half-line boundary density of the pentadiagonal
Toeplitz operator with symbol $f(\theta) = 2\cos\theta + 2\cos 2\theta$
at $\earw = e_{1} + e_{2}$ is $(\sin\theta + \sin 2\theta)^{2}/\pi$, by
extrapolation from the tridiagonal case where the boundary density at
$e_{1}$ is $\sin^{2}\theta / \pi$. This naive density fails the
Plancherel-norm test by a factor of $100+$: direct projection of $\earw$
onto the eigenbasis of $\Adj(L_{n-1})$ at $n = 200$ gives entries that
differ from $(\sin\theta_{j} + \sin 2\theta_{j})^{2} \cdot 2/(n+1)$ by
factors of order $10^{2}$. The correct density is the \emph{signed
angle gap}
\[
\rho_{w}(\lambda) \;=\; \frac{1}{\pi}\sin(\theta_{2}(\lambda) - \theta_{1}(\lambda))
\]
on the negative-spectrum interval $(-9/4, 0)$, where
$\theta_{1} < \theta_{2}$ are the two preimages of $\lambda$ under
$f$ (\Cref{thm:2path-stieltjes}). The naive form fails because the
eigenvectors of pentadiagonal Toeplitz are \emph{not} simple sinusoids
— a tridiagonal-specific structure that does not survive bandwidth $2$.

\subsection{F5. The asymmetric ``selector + dual'' framework}\label{sec:fm:asymmetric-dual}

A natural reflex after the universal-lemma failure (F1) is to combine
the max-degsum selector for $\dminus$ with a dual selector for $\dplus$,
hoping the asymmetry between $\splus$ and $\sminus$
(Elphick--Linz~\cite{ElphickLinz2023}) yields a tighter ceiling on the
positive side. Direct computational test on the same corpus rules this
out: positive-side and negative-side Rayleigh-tightness ratios for
Lemma~B1 (and its positive-side dual) agree to within $3\%$, with no
asymmetry to exploit. Moreover, the trace identity
$\dplus(v) + \dminus(v) = 4$ at any degree-$2$ ear forces a one-to-one
mapping between the two gains, so any one-sided selector for $\dminus$
is automatically a selector for $\dplus$ at the reflected threshold;
the dual framework collapses to a relabelling.

\subsection{F6. The $\alpha\omega$ fallback route}\label{sec:fm:alpha-omega}

Theorem~8.1 of \cite{AKMPZ2025} proves \Cref{conj:92} strictly under
the hypothesis $\alpha(\graph) \omega(\graph) \le n/17$. This route is
uniformly inapplicable to the $2$-tree class: every $2$-tree has
$\omega(\graph) = 3$ (treewidth $2$ caps the max clique at a triangle),
and the independence number $\alpha(\graph)$ scales linearly with $n$
on every binding family — empirically on the corpus
$(\alpha \omega)/n$ ranges between $0.8$ (books) and $1.0$ (paths),
\emph{never} below $1/17 \approx 0.059$. The constant $1/17$ in
Theorem~8.1 traces back to the slack $\epsilon = 1/16$ of the
$P_{3}$-removal lemma; improving $17$ requires sharpening that slack,
which is research upstream of our workstream. The $\alpha\omega$-route
is therefore not the right fallback for $2$-trees.

\subsection{F7. Sign error in Lemma~B1 optimisation}\label{sec:fm:sign-error}

An intermediate draft of \Cref{lem:B1} had the wrong sign in the
Rayleigh-quotient stationary equation
($\beta^{2} + (\Monem/\Wm)\beta - \Wm = 0$ instead of the correct
$\beta^{2} - (\Monem/\Wm)\beta - \Wm = 0$). This propagated into the
identification of $\beta_{+}$ as the minimiser: the wrong-sign analysis
identifies $\beta_{+} = (|\Monem| + \sqrt{(\Monem)^{2} + 4(\Wm)^{3}})/(2\Wm)$,
whereas the correct minimiser is
$\beta_{+} = (\sqrt{(\Monem)^{2} + 4(\Wm)^{3}} - |\Monem|)/(2\Wm)$.
The displayed lemma bound happens to be unchanged (a luck-of-the-algebra
coincidence: $R(\beta_{+}) = -\Wm/\beta_{+}$ rationalises to the same
final formula in both cases), but the intermediate Case-B slot
decomposition built on the wrong root identification gave a false
$\dminus(\maxdeg) \ge 1$ claim. Spot-check at $L_{5}$ catches it:
$\lambda_{\min}(\Adj(L_{5})) = -1.618$ exactly, while the wrong-sign
formula predicts $+1.387$. The corrected proof in
\Cref{sec:moment-form:lemma-B1} uses the right root throughout, and
the falsified $\dminus(\maxdeg) \ge 1$ claim is downgraded to the
conditional \Cref{lem:bminor-amin} plus the empirical
\Cref{prop:bminor-census}.

\subsection{F8. ``Floating-point certified'' $\neq$ interval-arithmetic certified}
\label{sec:fm:fp-vs-interval}

The $\dminus(\Ln) \ge 17/16 + 1/4$ certificate of \Cref{thm:2path-DK}
uses \texttt{numpy.linalg.eigvalsh} with an \emph{a-posteriori}
Demmel--Kahan forward-error bound \cite[Thm.~5.5]{DemmelANLA}. The
resulting statement is engineering-rigorous \emph{conditional on} the
coarse \texttt{dsyevd} envelope $c \le 10^{4}$: under that envelope
the forward error on each computed eigenvalue is provably bounded by
$8.88 \cdot 10^{-9}$ at $n \le 1000$ and the propagated $\dminus$
forward error by $1.42 \cdot 10^{-4}$, far below the slack $0.257$;
the envelope itself is plausible (over an order of magnitude above any
published constant for backward-stable symmetric eigensolvers) but not
formally derived for the exact LAPACK \texttt{dsyevd} pipeline within
this paper, and the certificate is \emph{not} interval arithmetic. A truly interval-arithmetic certificate (via
\texttt{mpmath} at, say, $50$ decimal digits, or via Krawczyk-style
verified eigensolver) would be a strictly stronger guarantee. For our
purposes the a-posteriori bound suffices; the upgrade path is
mechanical.

\subsection{F9. Simple-loop hypothesis for BBG-type asymptotic constants}
\label{sec:fm:simple-loop}

Bogoya--B\"ottcher--Grudsky \cite{BBG2015} and the related
Avram--Parter / Widom-type effective constants for finite-$n$
corrections to Szeg\H{o}'s theorem assume a \emph{simple-loop} symbol:
$f$ traces the unit circle once as $\theta$ traverses $[-\pi, \pi]$,
without singularities. Our pentadiagonal symbol
$f(\theta) = 2\cos\theta + 2\cos 2\theta$ has zeros at both
$\theta = \pi/3$ (interior, transversal) and $\theta = \pi$ (boundary),
so the simple-loop hypothesis fails. The \emph{rate} of convergence in
\Cref{thm:2path-stieltjes} — i.e.\ an explicit $n_{0}$ from which
$|\Iv - \Iinf{L}| \le \epsilon$ for given $\epsilon$ — is therefore not
available by off-the-shelf citation. The limit value
$\Iinf{L} \approx 1.0157$ and its closed form are unaffected; the rate
is recorded as open subobligation O13.1 of the workstream.

\subsection{F10. The $2$-path is not the universal binding case for
\texorpdfstring{$\Iv$}{I(v)}}\label{sec:fm:not-binding}

An intermediate working hypothesis held that the $2$-path family
asymptotic $\Iinf{L} \approx 1.0157$ is the binding case for the joint
invariant $\Iv$ at the max-degsum ear of every $2$-tree, i.e.\ that
$\Ivs \ge \Iinf{L}$ would close condition~(a) on general $2$-trees.
This is false at finite~$n$: at $n = 10$ the caterpillar $2$-tree
with graph6 \texttt{I\}qcHG\textasciigrave GO} satisfies
$\Ivs = 0.6384$, which is below $\Iinf{L}$ by $\approx 0.38$ and below
$\Iv$ on $L_{n}$ for every $n \in \{4, \ldots, 30\}$.

Two consequences. First, the open problem \Cref{prob:general-a} is
correctly framed as $\Ivs \ge \threshold = 0.4122$, not as
$\Ivs \ge \Iinf{L}$; the corpus minimum gives slack $\approx 0.23$
above $\threshold$ at the worst case. Second, a $2$-path-embedding
monotonicity argument cannot close condition~(a) on general $2$-trees,
since the caterpillar minimiser embeds a $2$-path that has
$\Iv > 0.6384$ on its own (the $2$-path on $5$ vertices already has
$\Iv = 0.79\ldots$). The structural lower bound \Cref{lem:cs-two-sided}
— a moment-form functional rather than a subfamily comparison — is
therefore the right tool. The caterpillar minimiser is preserved as a
permanent regression fixture in
\texttt{tests/fixtures/joint\_invariant\_falsified.json}.

\subsection{Summary table}\label{sec:fm:summary}

\Cref{app:fixtures} contains a tabular summary of failure modes
F1--F10, with columns: identifier, the precise false statement,
the smallest falsifying witness (combinatorial or numerical), and the
regression fixture file.
